%! Author = chtompki
%! Date = 1/14/26
% Example LaTeX document for GP111 - note % sign indicates a comment
\documentclass[12pt]{amsart}
% Default margins are too wide all the way around. I reset them here
\setlength{\hoffset}{-.75in}
\setlength{\textwidth}{6.5in}
\setlength{\voffset}{-.5in}
\setlength{\textheight}{9.0in}
\usepackage{amsmath}
\usepackage{amssymb}
\usepackage{amsthm}
\usepackage{pgfplots}
\usepackage{enumerate}
\usepackage{tikz}
\usetikzlibrary{shadows,arrows,arrows.meta,shapes,positioning,calc}% Required for the Circle and Triangle arrow tips
\usepackage{setspace}
\usepackage{siunitx}
\usepackage{enumitem}

\theoremstyle{conjecture}
\newtheorem{conjecture}{Conjecture}

\theoremstyle{definition}
\newtheorem{definition}{Definition}

\theoremstyle{question}
\newtheorem{question}{Question}

\theoremstyle{theorem}
\newtheorem{theorem}{Theorem}

\newenvironment{solution}
{\begin{proof}[Solution]}
{\end{proof}}

\title{Graph Theory:\\Homework 3\\ Book Section 1.3.}
\author{Rob Tompkins}
\date{\today}


\begin{document}

    \maketitle
    \noindent\textbf{1.} Determine the maximum number of edges in a bipartite subgraph of $P_n$, $C_n$,
    and $K_n$. (Hint: for $K_n$ it may be helpful to review the exercise from \S 1.1 on counting edges
    of bipartite graphs.)

    \begin{solution}
        Since $P_n$ is bipartite, all of its $n-1$ edges may be retained. Thus the maximum is $n-1$.

        If $n$ is even, then $C_n$ is bipartite, so the maximum is $n$. If $n$ is odd, then $C_n$ is not
        bipartite, so at least one edge must be omitted. Deleting any one edge leaves a path with $n-1$
        edges, which is bipartite. Hence the maximum is $n-1$ when $n$ is odd.

        Finally, let a bipartite subgraph of $K_n$ have partite sets of sizes $a$ and $b$. We may include
        isolated vertices in the two sets, so we can assume $a+b=n$. Every edge goes between the two sets,
        and hence there are at most
        \[
            ab\leq \left\lfloor\frac{n^2}{4}\right\rfloor
        \]
        edges; the product is largest when the two integers are as equal as possible. This bound is attained
        by the subgraph $K_{\lfloor n/2\rfloor,\lceil n/2\rceil}$. Therefore the desired maximum for $K_n$ is
        $\lfloor n^2/4\rfloor$.
    \end{solution}

    \noindent\textbf{2.} Prove that an even graph has no cut-edge

    \begin{proof}
        Suppose, to the contrary, that an even graph $G$ has a cut-edge $e=uv$. In $G-e$, let $H$ be the
        component containing $u$. Every vertex of $H$ other than $u$ has the same (even) degree in $H$ as in
        $G$, while $u$ has degree $d_G(u)-1$, which is odd. Thus $H$ has exactly one vertex of odd degree.
        This contradicts the Handshaking Lemma, according to which every graph has an even number of
        odd-degree vertices. Hence an even graph has no cut-edge.
    \end{proof}

    \noindent\textbf{3.} The problem involves the hypercube $Q_k$.
    \begin{itemize}
        \item[(a)] Use the recursive description of $Q_k$ (Example 1.3.8) to prove that
        $|e(Q_k)| = k2^{k-1}$
        \item[(b)] Prove that $K_2,3$ is not contained in any hypercube $Q_k$.
        (Hint: prove that two distinct vertices in $Q_k$ cannot have three common neighbors.)
    \end{itemize}

    \begin{proof}
        \begin{enumerate}[label=(\alph*)]
            \item The recursive construction of $Q_k$ takes two copies of $Q_{k-1}$ and adds a matching
            joining corresponding vertices. Since $Q_{k-1}$ has $2^{k-1}$ vertices, this gives
            \[
                |E(Q_k)|=2|E(Q_{k-1})|+2^{k-1}.
            \]
            The formula is true for $Q_1$, which has one edge. Assuming
            $|E(Q_{k-1})|=(k-1)2^{k-2}$, the recurrence yields
            \[
                |E(Q_k)|=2(k-1)2^{k-2}+2^{k-1}=k2^{k-1}.
            \]
            Thus the result follows by induction.

            \item View the vertices of $Q_k$ as binary strings of length $k$, with adjacency when two
            strings differ in exactly one coordinate. Let $x$ and $y$ be distinct vertices having a common
            neighbor $z$. Since both $x$ and $y$ differ from $z$ in one coordinate, $x$ and $y$ must differ
            in exactly two coordinates. A common neighbor is obtained by changing in $x$ one of those two
            coordinates, so there are exactly two possible common neighbors. (If $x$ and $y$ do not differ
            in exactly two coordinates, they have none.) Consequently two distinct vertices of $Q_k$ can
            never have three common neighbors. In a copy of $K_{2,3}$, however, the two vertices in the
            part of size two have all three vertices of the other part as common neighbors. Therefore no
            $Q_k$ contains $K_{2,3}$.
        \end{enumerate}
    \end{proof}

    \noindent\textbf{4.} Prove that every edge in the Petersen graph $G$ belongs to exactly four 5-cycles,
    and use this to show that the Petersen graph has exactly twelve 5-cycles.
    (Hint: For the first part, extend the edge to a copy of $P_4$, and apply Proposition 1.1.38.)

    \begin{proof}
        Use the standard realization of the Petersen graph whose vertices are the $2$-subsets of
        $\{1,2,3,4,5\}$, with two vertices adjacent when the subsets are disjoint. Permutations of
        $\{1,2,3,4,5\}$ induce automorphisms, and any edge can therefore be carried to the edge
        $\{1,2\}\{3,4\}$. It is enough to count the 5-cycles through this edge.

        The neighbors of $\{3,4\}$ other than $\{1,2\}$ are $\{1,5\}$ and $\{2,5\}$, while the neighbors
        of $\{1,2\}$ other than $\{3,4\}$ are $\{3,5\}$ and $\{4,5\}$. Choosing one vertex from each of
        these two pairs determines the middle vertex of a path back from $\{3,4\}$ to $\{1,2\}$: it must
        be the unique $2$-subset disjoint from both chosen vertices. The four choices give middle vertices
        $\{2,4\},\{2,3\},\{1,4\},\{1,3\}$, respectively. Thus there are exactly four 5-cycles through the
        fixed edge, and hence through every edge.

        Now count incidences $(e,C)$ in which $e$ is an edge lying on a 5-cycle $C$. The Petersen graph is
        $3$-regular on $10$ vertices, so it has $15$ edges. Counting first by edges gives $15\cdot4=60$
        incidences. If there are $N$ 5-cycles, counting by cycles gives $5N$ incidences. Therefore
        $5N=60$, and the Petersen graph has $N=12$ 5-cycles.
    \end{proof}

    \noindent\textbf{5.} Use complete graphs and counting arguments (not algebra!) to prove that
    \begin{itemize}
        \item[(a)] $\binom{n}{2} = \binom{k}{2} + k(n - k) + \binom{n-k}{2}$ for all integers $0\le k\le n$.
        (Hint: Count the edges of the complete graph $K_n$ in two different ways.)
        \item[(b)] If $\sum n_i = n$ then $\sum\binom{n_i}{2} \le \binom{n}{2}$
    \end{itemize}

    \begin{proof}
        \begin{enumerate}[label=(\alph*)]
            \item Partition the vertices of $K_n$ into sets $A$ and $B$ of sizes $k$ and $n-k$.
            There are $\binom{k}{2}$ edges with both endpoints in $A$, $k(n-k)$ edges with one endpoint
            in each set, and $\binom{n-k}{2}$ edges with both endpoints in $B$. These three disjoint classes
            contain every edge of $K_n$. Since $K_n$ has $\binom{n}{2}$ edges, the identity follows.

            \item Partition the $n$ vertices of $K_n$ into sets of sizes $n_1,n_2,\ldots$. Within the
            $i$th set there are $\binom{n_i}{2}$ edges, so $\sum_i\binom{n_i}{2}$ counts all edges whose
            endpoints lie in the same part. These form only a subset of the $\binom{n}{2}$ edges of $K_n$;
            the remaining edges join distinct parts. Hence
            $\sum_i\binom{n_i}{2}\leq\binom{n}{2}$.
        \end{enumerate}
    \end{proof}


\end{document}